\section{Logical meaning and exact polynomial definitions}
\label{sec:foundation}
Throughout, $\N=\{0,1,2,\ldots\}$, $\Z_{>0}=\{1,2,3,\ldots\}$ and
$\Z$ is unrestricted. The domain belongs to an equation's definition.
An equation over positive integers is not an unrestricted-integer equation;
a natural-domain chart needs a separate sign-restoration proof.

Let $\operatorname{Ref}_{\mathrm{ZFC}}(r)$ be the fixed arithmetized
predicate saying that $r$ codes a finite syntactic ZFC refutation. For
each selected polynomial $P$ on its stated domain $D$, the ordinary claim is
\begin{equation}\label{eq:calibration}
 \exists x\in D^n\;P(x)=0
 \quad\Longleftrightarrow\quad
 \exists r\in\N\;\operatorname{Ref}_{\mathrm{ZFC}}(r).
\end{equation}
The unbounded arithmetic transformations below compose with the explicitly
identified inherited calibration. Formalizing this proof chain in ZFC yields
\begin{equation}\label{eq:consistency}
 \mathrm{ZFC}\vdash
 \bigl[\forall x\in D^n\;P(x)\ne0\bigr]
 \longleftrightarrow\Con(\mathrm{ZFC}).
\end{equation}
This is an equivalence, not a proof of either side. It does not assert
consistency, decide solvability, give a globally smallest polynomial, or
provide a new end-to-end proof-assistant kernel check.

\subsection{The retained calibration boundary}
For the direct arithmetic families, the named logical interfaces are
\begin{quote}\small
\file{KEndpoint.c222_iff_zf_refutation}\\
\file{ConstructibleChoice.zf_refutation_iff_zfc_refutation}.
\end{quote}
The first relates the fixed numerical calculus to a finite ZF refutation;
the second transports between ZF and ZFC refutations. Their exact sources,
immutable historical pins and ordinary arithmetic bridges are retained in
the combined source closure. The coefficient-height-one families have
separately identified original quadratic sources and transfer their inherited
zero-set calibrations through the explicit graph maps below. They are not
silently rebased to the newest direct system.

The positive scalar branch uses its own machine, native-encoding and
positive-P66 calibration at one complete tuple of fixed coefficients.
Its source theorem is not inferred merely from the sizes of those numerals.
The whole tuple, exact finite language, coefficient recipes and unbounded
native proof are supplied; replacing a single port in isolation is not
licensed. The source package's old pin and this round's current comparison
pin are different records and are never substituted for each other.

There are four different kinds of evidence here. The logical calibration
relates finite syntax to the original arithmetic source. Unbounded proofs
then establish Pell classification, cofinal witness construction, CRT
compatibility and sign recovery. Exact algebra transfers the source zero
set to finite new arrays or circuits. Finally independent finite checkers
reconstruct those actual arrays, not merely an abstract theorem instance.
A hash proves byte identity, not a theorem. Coefficient reconstruction proves
an identity, not an unbounded Pell assertion. Successful finite replay is
not an end-to-end formal derivation of~\eqref{eq:calibration}.

\subsection{Literal arrays and graphs}
For each residual endpoint, the ordered \texttt{variables} array in
\file{equation.json} fixes coordinates $x_0,\ldots,x_{n-1}$. A term
$[c,[i_1,\ldots,i_d]]$ denotes $c\prod_jx_{i_j}$; an empty index list
denotes the constant $c$. Repeated indices mean ordinary powers. If the
listed residuals are $R_0,\ldots,R_{m-1}$, the literal polynomial is
\begin{equation}\label{eq:sos}
 P(x)=\sum_{j=0}^{m-1}R_j(x)^2.
\end{equation}
The sum of squares is zero exactly when every row is zero, on each of the
stated integer domains. The degree of a merged reciprocal row need not be
two; the advertised total degree is separately checked after collection.
The quartic endpoints have degree-at-most-two rows and a nonzero collected
quartic coefficient. Adjacent collected arrays give every actual integer
coefficient, including a nonzero constant. There is no variable exponent,
factorial, division, radical or sequence primitive in those arrays.

The height-one endpoints are defined by their complete collected
coefficient arrays. Their graph certificates establish a nonnegative
energy identity with an explicit integral lift and inverse projection.
An independent checker reconstructs every energy coefficient and compares
the complete output, including equality and unit terms.

A scalar circuit is a finite directed acyclic graph whose leaves are the
stated independent variables or supplied fixed integer constants. Every
internal gate is binary addition, subtraction or multiplication. It therefore
denotes a unique ordinary polynomial over $\Z$. Registers are abbreviations,
not new quantified unknowns. Fixed powers in the printed formula stand for
the specified multiplication gates; no quotient or square-root gate is used.
All numerical ports are supplied in full, not replaced by a hash or a
promise that a suitable large integer exists.

\subsection{Metrics and finite comparisons}
For a nonzero collected polynomial $P=\sum_\alpha c_\alpha x^\alpha$,
\[
 \deg P=\max_{c_\alpha\ne0}|\alpha|,\qquad
 |\supp P|=\#\{\alpha:c_\alpha\ne0\},\qquad
 \htpoly(P)=\max_\alpha|c_\alpha|.
\]
Support includes the constant if nonzero. Residual count is not support.
Every independent coordinate must be live after collection. Scalar cost
counts live binary arithmetic gates, including squaring and multiplication
by fixed constants; reading a leaf and copying a register are free. Dead
gates do not count. Circuit gate count is not bit complexity. When no full
scalar expansion is supplied, its support and full coefficient height are
unmeasured, not zero and not the maximum literal leaf.

Pareto comparisons hold only among the listed, certified presentations on
compatible domains, with degree fixed where stated. A point dominates
another if it is no worse in every compared measure and strictly better in
at least one. Smaller dimension, lower degree, sparser support, coefficient
height one and fewer circuit gates are distinct objectives. No worldwide
priority or global minimality conclusion follows from this finite catalog.
